% Vietnamese translation for OI Public Library
% Source-File: IOI中国国家候选队论文/2007/day1/5.杨沐《浅析信息学中的“分”与“合”》.ppt
\documentclass[11pt]{article}
\usepackage{oipl-vn}

\newcommand{\Oh}{\mathcal{O}}
\newcommand{\len}{\operatorname{len}}

\title{Bản trình chiếu: Tư duy phân và hợp trong tin học}
\author{Yang Mu}
\date{2007}

\begin{document}
\maketitle

\origtitle{Slide companion for division and combination in informatics}
\sourcepath{IOI China candidate team papers / 2007 / day 1 / Yang Mu.ppt}

\section{Phạm vi bản dịch}

Tệp PowerPoint đi cùng bài viết đã được dịch từ PDF. Bản ghi chú này giữ lại
tinh thần của slide: dùng việc tách nhỏ và gom lại như một phương pháp suy
nghĩ, không chỉ như thuật toán chia để trị theo nghĩa hẹp.

\section{Tư duy phân}

Khi bài toán quá lớn hoặc có quá nhiều trường hợp, ta thêm một tham số, một
điều kiện phụ hoặc một cách phân lớp để làm cấu trúc rõ hơn. Trong ví dụ
\textit{Milk Patterns}, thay vì tìm trực tiếp đoạn dài nhất xuất hiện ít nhất
\(k\) lần, ta hỏi với một độ dài \(\len\) cố định liệu có đoạn nào xuất hiện
đủ nhiều lần hay không. Tính đơn điệu theo \(\len\) dẫn tới tìm kiếm nhị phân.

\section{Tư duy hợp}

Chiều ngược lại là gom các phần đã hiểu thành một cấu trúc lớn hơn. Khi các
đối tượng nhỏ có quan hệ qua khoảng cách, lớp dư hoặc cạnh trong đồ thị, việc
nhìn chúng như một hệ thống chung có thể làm lộ ra bài toán đường đi ngắn nhất,
quy hoạch động hoặc dựng cây.

\section{Ví dụ kỹ thuật}

Trong các slide có xuất hiện những khuôn mẫu như phân lớp theo modulo \(K\),
đồ thị \(G=(V,A,C)\), chạy SPFA trên trạng thái đã tách, và các công thức
quy hoạch động dạng \(DP[i,j]\). Điểm chung là mỗi lần tách đều phải phục vụ
một bước hợp lại rõ ràng; nếu không, số trường hợp chỉ tăng mà không giúp
thuật toán đơn giản hơn.

\section{Kết luận}

``Phân'' và ``hợp'' là hai thao tác bổ sung cho nhau. Tách đúng giúp bài toán
có tính đơn điệu, độc lập hoặc truy hồi; hợp đúng giúp tổng hợp lời giải của
các phần mà không đánh mất ràng buộc toàn cục. Đây là kỹ năng thiết kế thuật
toán hơn là một công thức cố định.

\end{document}
